%%%%%%%%%%%%%%%%%%%%%%%%%%%%%%%%%%%%%%%%%%%%%%%%%%%%%%%
% % Lines starting with % are comments, which are ignored.
% % This is a handy way of indicating the date and version of
% % your document, to wit:
% %
% % LaTeX sample file
% % Modified March, 2002
% %
%%%%%%%%%%%%%%%%%%%%%%%%%%%%%%%%%%%%%%%%%%%%%%%%%%%%%%%
\DocumentMetadata{
  % PDF standards: ONE key whose value is a comma-separated list --
  % UA-2 (accessibility) plus A-4f (archival; the archival companion of
  % UA-2). Do NOT use separate pdfstandard=... keys: with ordinary
  % key-value parsing only the last one takes effect.
  pdfstandard={UA-2,A-4f},
  lang=en-US,
  tagging=on,
  tagging-setup={
    math/setup=mathml-SE
  }
}

\documentclass{ltx-talk}
% lecture-macros.tex -- shared package setup for lecture-template.tex
% (ltx-talk slides / handouts).
%
% NOTE ON PATHS: every template in this repository is stored one level up
% from where it is compiled (see README.md). All relative \input paths
% below are resolved against the folder the document is compiled in,
% which sits one directory above this file.

% --- Tables, bookmarks, strikethrough ------------------------------------
% booktabs: professional table rules (\toprule, \midrule, \bottomrule)
% bookmark: PDF outline entries (PDF bookmarks)
% ulem[normalem]: \sout strikethrough that survives math mode
\usepackage{booktabs}
\usepackage{bookmark}
\usepackage[normalem]{ulem}
\usepackage{multicol}

% --- Math ------------------------------------------------------------
% amsmath: equation environments and numbering
% mathrsfs: \mathfrak  |  stmaryrd: extra symbols
\usepackage{amsmath}
\usepackage{mathrsfs}
\usepackage{stmaryrd}

% --- Layout -------------------------------------------------------------
% float: [H] float placement
% fontspec: Unicode/OpenType fonts (compile with XeLaTeX or LuaLaTeX)
\usepackage{float}
\usepackage{fontspec}

% --- Colors ---------------------------------------------------------------
% dvipsnames: extended named color palette (the Accessible* colors below
% are built on it);
% table: \rowcolors / \cellcolor for shading tables.
% NOTE: xcolor may only be loaded ONCE per document -- a second
% \usepackage{xcolor} with different options is a fatal "Option clash"
% error, so ALL of its options are listed here.
\usepackage[dvipsnames,table]{xcolor}

% --- Multimedia -----------------------------------------------------------
% animate: embeds audio/video (and legacy Flash) media in a frame.
\usepackage{animate}

% --- Bibliography ----------------------------------------------------------
% biblatex (biber backend): citations for lecture notes.
\usepackage[
  backend=biber,
  style=alphabetic,
  url=true,
  doi=true,
  isbn=true,
  backref=true
]{biblatex}

% --- Colored boxes ----------------------------------------------------------
% talkthemetcolorbox: the colored-box theme for ltx-talk slides
% (shadow/rounded/shaded blocks, showtitle, titlepage styles).
\usepackage[
  shadow,
  rounded,
  shaded,
  blocks,
  showtitle,
  titlepage
]{talkthemetcolorbox}

% --- "theindex" helper environment ------------------------------------------
% Provides a lightweight theindex environment (an itemize) and the
% \@idxitem/\subitem/\subsubitem/\indexspace macros that LaTeX's standard
% theindex machinery expects, so hand-written index pages indent correctly.
\makeatletter
\newenvironment{theindex}
  {\begin{itemize}}
  {\end{itemize}}
\providecommand\@idxitem{\par\hangindent 40\p@}
\providecommand\subitem{\@idxitem \hspace*{20\p@}}
\providecommand\subsubitem{\@idxitem \hspace*{30\p@}}
\providecommand\indexspace{\par \vskip 10\p@ \@plus5\p@ \@minus3\p@\relax}
\makeatother

% --- Index -----------------------------------------------------------------
% imakeidx: builds the printed index.
% The redefinition below appends "|hyperpage" to every index entry so that,
% in the PDF, each index entry is a link to its page.
\usepackage{imakeidx}
\makeatletter
\let\imki@orig@wrindexentrysplit\imki@wrindexentrysplit
\renewcommand\imki@wrindexentrysplit[3]{%
  \imki@orig@wrindexentrysplit{#1}{#2|hyperpage}{#3}%
}

\makeatletter
% --- Section titles: capture + one bookmark per logical section -------------
% \section is redefined so that:
%   1. the section title is captured into the plain-TeX register
%      \mytalksectitletoks (reprinted verbatim on the auto divider frame,
%      see the cmd/section/after hook below -- plain-TeX storage is
%      catcode-safe, unlike expanding the title in a hook), and
%   2. exactly ONE depth-1 PDF bookmark is created per \section call
%      (hyperref's own bookmark is suppressed via \Hy@writebookmark).
\DeclareCommandCopy{\OriginalSectionForBookmarks}{\section}
\newtoks\mytalksectitletoks

\RenewDocumentCommand{\section}{s o m}{%
  \begingroup
    \def\Hy@writebookmark##1##2##3##4##5{}%
    \global\mytalksectitletoks={#3}%
    \IfBooleanTF{#1}
      {\phantomsection\IfNoValueTF{#2}{\OriginalSectionForBookmarks*{#3}}{\OriginalSectionForBookmarks*[#2]{#3}}}%
      {\IfNoValueTF{#2}{\OriginalSectionForBookmarks{#3}}{\OriginalSectionForBookmarks[#2]{#3}}}%
  \endgroup
  \IfBooleanTF{#1}
    {\IfNoValueTF{#2}{\pdfbookmark[1]{#3}{section-star.\thepage.\number\inputlineno}}{\pdfbookmark[1]{#2}{section-star.\thepage.\number\inputlineno}}}%
    {\IfNoValueTF{#2}{\pdfbookmark[1]{Section \thesection: #3}{section.\thesection}}{\pdfbookmark[1]{Section \thesection: #2}{section.\thesection}}}%
}
\makeatother

% --- One PDF bookmark per logical frame (NOT per overlay) ---------------------
% ltx-talk writes one PDF page per overlay, so a naive frametitle hook would
% register a new bookmark for every overlay of the same frame. The
% \g_slideutils_frame_bookmark_int register remembers the last frame that
% got a bookmark, so subsequent overlays of the same frame are skipped.
% This is also what keeps BeamerPresenter's cumulative annotations working:
% one logical frame must stay one logical unit in the PDF outline.
\ExplSyntaxOn
\int_new:N \g_slideutils_frame_bookmark_int

\AddToHook{cmd/frametitle/after}{%
  \tl_if_empty:NF \g__talk_frame_title_tl
    {
      \ifnum\value{frame}=\g_slideutils_frame_bookmark_int\relax
      \else
        \int_gset:Nn \g_slideutils_frame_bookmark_int { \value{frame} }%
        \pdfbookmark[2]{\tl_use:N \g__talk_frame_title_tl}
          {frame.\int_use:N \g_slideutils_frame_bookmark_int}
      \fi
    }
}
\ExplSyntaxOff
\makeatother

% --- Printed index + theme for the colored boxes -------------------------------
% \makeindex[...]: two-column index titled "Index", listed in the TOC.
% The two EditInstance calls style the standard header/footer instances
% of talkthemetcolorbox (cyan background; footer order title/author/
% date/framenumber).
\makeindex[columns=2, title=Index, intoc]\EditInstance{header}{std}{
  background-color = cyan,
  color            = black,
  height           = 0.75cm
}
\EditInstance{footer}{std}{
  background-color = cyan,
  element-order = {title, author, date, framenumber},
}

% --- Auto section divider frame ---------------------------------------------------
% After each \section, insert a centered, large, structure-colored frame
% that prints the section title (from \mytalksectitletoks).
\AddToHook{cmd/section/after}{%
  \begin{frame}
    \centering
    \color{structure}
    \LARGE
    \the\mytalksectitletoks
  \end{frame}%
}

% --- Tagging sockets (tagged PDF structure) ---------------------------------------
% ltx-talk exposes tagging "sockets" where the PDF structure can attach
% section titles; register one and leave it unassigned (none) so the
% class's own tagging keeps control.
\ExplSyntaxOn
\AtBeginDocument{
  \NewTaggingSocketPlug { talk / sec / title } { none } { }
  \AssignTaggingSocketPlug { talk / sec / title } { none }
}
\ExplSyntaxOff

% --- Accessibility roles ----------------------------------------------------------------
% Map \frametitle to the H1 heading role in the tagged PDF, so screen
% readers treat slide titles as headings.
\tagpdfsetup{role / new-tag = frametitle / H1}

\makeatletter
% --- Page label = frame number -------------------------------------------------------
% Redefine what feeds the PDF page label to track the frame counter,
% not the physical page counter. Replace \theframenumber below with
% whatever internal macro/property ltx-talk actually uses to produce
% the "framenumber" footer element — confirm via ltx-talk-frame.dtx.
\renewcommand\thepage{\arabic{frame}}
\makeatother

\addbibresource{refs.bib}

% --- "Problem Time!" slides ------------------------------------------------------------------
% \ProblemSlide{problem}{solution} emits two frames: the problem, then
% its solution after a \pause, each numbered by the problem counter.
\newcounter{problemcounter}
\setcounter{problemcounter}{0}
\newcommand{\ProblemSlide}[2]{
  \stepcounter{problemcounter}
  \begin{frame}
    \frametitle{Problem Time!}
    Problem \theproblemcounter: #1
  \end{frame}

  \begin{frame}
    \frametitle{Problem Time!}
    Solution to Problem \theproblemcounter: \pause

    #2
  \end{frame}
}
% \highlight[<colour>]{<stuff>}
% from https://tex.stackexchange.com/questions/33401/a-version-of-colorbox-that-works-inside-math-environments
\newcommand{\highlight}[2][yellow]{\mathchoice%
  {\colorbox{#1}{$\displaystyle#2$}}%
  {\colorbox{#1}{$\textstyle#2$}}%
  {\colorbox{#1}{$\scriptstyle#2$}}%
  {\colorbox{#1}{$\scriptscriptstyle#2$}}}%

% --- Fonts ------------------------------------------------------------------------------
% Lato (text) + LeteSansMath (math): open, accessible typefaces.
\setmainfont{Lato}[
  Ligatures=TeX,
  Numbers={Lining,Monospaced}
]
\setsansfont{Lato}[
  Ligatures=TeX,
  Numbers={Lining,Monospaced}
]
\setmathfont{LeteSansMath.otf}\usepackage{emoji}

% --- Accessible colors ---------------------------------------------------------------------
% WCAG-contrast-safe variants of red/green for text on white:
% red!80!black and green!40!black both clear the 4.5:1 requirement.
% AltYellow: soft background for highlight boxes.
\colorlet{AccessibleRed}{red!80!black}
\colorlet{AccessibleGreen}{green!40!black}
\definecolor{AltYellow}{RGB}{255,245,170}

% --- Links ---------------------------------------------------------------------------------
% hyperref + link colors: every link uses AccessibleGreen (7.2:1 on
% white), so internal, citation, and URL links all pass WCAG 1.4.3.
\usepackage{hyperref}
\hypersetup{
  colorlinks=true,
  linkcolor=AccessibleGreen,
  urlcolor=AccessibleGreen,
  citecolor=AccessibleGreen,
}

% --- Shared course macros ---------------------------------------------------------------------
% Root macros.tex: math shorthands used across all course documents
% (\RR, \NN, \ZZ, \QQ, \grad, \curl, \tr, \der, \pder, ...).
\newcommand{\RR}{\mathbb{R}}
\newcommand{\NN}{\mathbb{N}}
\newcommand{\ZZ}{\mathbb{Z}}
\newcommand{\QQ}{\mathbb{Q}}
\newcommand{\vv}{\vec{v}}
\newcommand{\ww}{\vec{w}}
\newcommand{\xx}{\vec{x}}
\newcommand{\yy}{\vec{y}}
\newcommand{\zz}{\vec{z}}
\newcommand{\uu}{\vec{u}}
\newcommand{\basis}[1]{\vec{e_{#1}}}
\newcommand{\id}{\mathrm{id}}
\newcommand{\CC}{\mathbb{C}}
\renewcommand{\div}{\operatorname{div}}
\newcommand{\grad}{\operatorname{grad}}
\newcommand{\curl}{\operatorname{curl}}
\newcommand{\tr}{\operatorname{tr}}
\newcommand{\der}[2]{\dfrac{d #1}{d #2}}
\newcommand{\pder}[3][]{\dfrac{\partial^{#1} #2}{\partial #3}}


% % Title and author(s)
%%%%%%%%%%%%%%%%%%%%%%%%%%%%%%%%%%%%%%%%%%%%%%%%%%%%%%%
\title{Lecture 12 - Inverse and Implicit Differentiation}
\author{Garrett Credi}
\date{10/20/2026}
%%%%%%%%%%%%%%%%%%%%%%%%%%%%%%%%%%%%%%%%%%%%%%%%%%%%%%%
% %
% % The next command allows your in import encapsulated
% % postscript files, .epsf or .eps files, which
% % contain vector graphic image data.
% %
%%%%%%%%%%%%%%%%%%%%%%%%%%%%%%%%%%%%%%%%%%%%%%%%%%%%%%%
% %   The first thanks indicates your affiliation
% %
% %  Just the name here.
% %
% % Your mailing address goes at the end.
% %
% % \thanks is also how you indicate grant support
% %
%%%%%%%%%%%%%%%%%%%%%%%%%%%%%%%%%%%%%%%%%%%%%%%%%%%%%%%

%%%%%%%%%%%%%%%%%%%%%%%%%%%%%%%%%%%%%%%%%%%%%%%%%%%%%%%
% $ The next command ensures the margins are 1 inch on
% % all sides.
% %
%%%%%%%%%%%%%%%%%%%%%%%%%%%%%%%%%%%%%%%%%%%%%%%%%%%%%%%
\begin{document}
\begin{frame}
  \maketitle
\end{frame}

\begin{frame}
  \frametitle{Updates}

  Curve information forthcoming - in contact with course coordinator about adjusting final grade distribution.

  Regrade requests for Exam 1 closed.
\end{frame}

\begin{frame}
  \frametitle{Quiz Time!}

  Let's get into neat rows, spread out, and slay this quiz!
\end{frame}

\section{Inverse and Implicit Differentiation}

\begin{frame}
  \frametitle{The Road So Far}
      Recall we're trying to build a \textit{synthetic} approach to studying derivatives: \pause

    Atomic Rules: \pause
  \begin{multicols}{2}
    \begin{itemize}
      \item Power Rule: $\dfrac{d}{dx} x^n$ 
      \item Trigonometric Rules: $\dfrac{d}{dx} \sin(x)$ etc.
      \item Exponential/Logarithm Rules: $\dfrac{d}{dx} e^x$, $\dfrac{d}{dx} \ln(x)$
    \end{itemize}
  \end{multicols}

  Combination Rules: \pause
  \begin{multicols}{2}
    \begin{itemize}
      \item Sum/Difference Rule: $\dfrac{d}{dx} (f(x) \pm g(x))$ 
      \item Product Rule: $\dfrac{d}{dx} f(x)g(x)$ 
      \item Chain Rule: $\dfrac{d}{dx} g(f(x))$ 
      \item Constant Multiple Rule: $\dfrac{d}{dx} c \cdot f(x)$
      \item Quotient Rule: $\dfrac{d}{dx} \dfrac{f(x)}{g(x)}$ 
      \item Inverse Function Rule: $\dfrac{d}{dx} f^{-1}(x)$
    \end{itemize}
  \end{multicols}
\end{frame}

\begin{frame}
  \frametitle{Explicit vs. Implicit Definitions I} 

  Up until now we've been dealing with \textit{explicit} definitions
  of various functions and have been using those \textit{explicit}
  definitions to compute their derivatives. \pause

  An \textit{explicit} definition of a function is one where
  we defined $f(x)$ by \textit{explicitly} saying where a given input $x$
  maps to. E.g.: \pause

  \begin{multicols}{3}
    \begin{itemize}
      \item $f(x) = x^2$ 
      \item $g(x) = \sin(x)$
      \item $h(x) = \sqrt{x^2 + 2x}$
      \item $j(x) = e^{\sin(x)}$
      \item $k(x) = (x + \sin(x))^{\frac{1}{2}}$
      \item etc
    \end{itemize}
  \end{multicols} \pause

  \textcolor{red}{This is not the only way in which we can define functions!}
\end{frame}

\begin{frame}
  \frametitle{Explicit vs Implicit Definitions II}

  We could instead define relations \textit{implicitly} - 
  i.e. by writing out some equation that we want the function
  to satisfy, e.g.: \pause

  \begin{multicols}{3}
    \begin{itemize}
      \item $x^2 + y^2 = 1$
      \item $y^2 = x^3 + 2x + 1$
      \item $\arctan(\frac{y}{x}) = \dfrac{\pi}{4}$
      \item $\sin(x) + \cos(y) = 1$
      \item $x^2 - y^2 = 2$
      \item $x^3 + y^3 = 1$
    \end{itemize}
  \end{multicols} \pause

  Just as we can \textit{graph} a function:
  i.e. we plot all the points $(x, f(x))$, \pause
  we can plot some implicit equation $F(x, y) = 0$
  by plotting all the points $(x, y)$ so that $F(x, y) = 0$. \pause

  When is this useful?
\end{frame}

\begin{frame}
  \frametitle{Applications of Implicit Relations}

  Oftentimes, real world process don't have simple ways of relating one quantity to another quantity. \pause

  Consider: conservation of momentum/energy in physics:
  \begin{gather}
    \sum\limits_{i = 1}^n m_i v^{init}_i = \sum\limits_{i = 1}^n m_i v^{final}_i \\
    \sum\limits_{i = 1}^n \dfrac{1}{2} m_i (v^{init}_i)^2 = \sum\limits_{i = 1}^n \dfrac{1}{2} m_i (v^{final}_i)^2
  \end{gather} \pause

  The relation between initial and final velocities here are not a simple
  input-to-output relationship - \pause they are equations that need to be enforced. \pause

  Also: predator-prey models, resource allocation, budgeting, etc.
\end{frame}

\begin{frame}
  \frametitle{Inverse Functions as a Relation}

  To address our final desired synthetic rule (inverse function rule) \pause
  we can also see that inverse functions are defined via an implicit rule: \pause

  \begin{block}{Inverse Function}
    For a function $f(x)$ with domain $D \subseteq \mathbb{R}$ and range $R \subseteq \mathbb{R}$, the inverse function $f^{-1}(x)$ is the function that: \pause
    \begin{itemize}
      \item Has domain $R$ (domain of inverse is range of original) \pause
      \item Has range $D$ (range of inverse is domain of original) \pause
      \item Satisfies the two relations: 
        \begin{enumerate}
          \item $f(f^{-1}(x)) = x$ 
          \item $f^{-1}(f(x)) = x$
        \end{enumerate} \pause
        i.e. $f(x)$ and $f^{-1}(x)$ ``undo" eachother.
    \end{itemize}
  \end{block}
\end{frame}

\begin{frame}
  \frametitle{Examples of Inverse Functions}
  For some examples: \pause
  \begin{enumerate}
    \item If $f(x) = x^2$ with domain \textit{the positive real numbers} and range \textit{the positive real numbers} \pause $f^{-1}(x) = \sqrt{x}$ is the positive square root. \pause
    \item If $f(x) = x^2$ with domain \textit{the negative real numbers} and range \textit{the positive real numbers} \pause $f^{-1}(x) = - \sqrt{x}$
      is the negative square root. \pause
    \item If $f(x) = \sin(x)$ with domain $[0, 2\pi)$ and range $[-1, 1]$ \pause $f^{-1}(x) = \arcsin(x)$  \pause
    \item If $f(x) = e^x$ with domain $\RR$ and range $(0, \infty)$ \pause
      $f^{-1}(x) = \ln(x)$.
  \end{enumerate}
\end{frame}

\begin{frame}
  \frametitle{Inverse Differentiation I}

  Question: Using the relation
  \begin{equation}
    f(f^{-1}(x)) = x
  \end{equation}
how might we relate the derivative of $f(x)$ with the derivative of it's inverse function $f^{-1}(x)$? \pause

Answer: Chain Rule! \pause

\begin{gather}
  f(f^{-1}(x)) = x \\
  \text{Differentiate both sides} \\ \pause
  \dfrac{d}{dx} f(f^{-1}(x)) = \dfrac{d}{dx} x \\ \pause
  \left(\dfrac{d}{dx} f^{-1}(x)\right) \cdot f'(f^{-1}(x)) = 1 \\ \pause
  \dfrac{d}{dx} f^{-1}(x) = \dfrac{1}{f'(f^{-1}(x))}
\end{gather} 
 \end{frame}

 \begin{frame}
   \frametitle{Inverse Differentiation II}

   \begin{alertblock}{Inverse Differentiation}\index{derivative!of an inverse function}
     If $f(x)$ and $f^{-1}(x)$ are an inverse pair of functions
     then: \pause
     \begin{equation}
       \dfrac{d}{dx} f^{-1}(x) = \dfrac{1}{f'(f^{-1}(x))}
     \end{equation}
   \end{alertblock} \pause
 \end{frame}
 \begin{frame}
   \frametitle{Inverse Differentiation III}

   Ex. $\dfrac{d}{dx}\ln(x)$. \pause

   $\ln(x)$ is inverse to $e^x$ $\implies$ $f(x) = e^x, f^{-1}(x) = \ln(x)$. \pause

   \begin{align}
     \dfrac{d}{dx} \ln(x) &= \dfrac{1}{\left(\dfrac{d}{du} e^u\right)\mid_{u = \ln(x)}} \pause = \dfrac{1}{\left(e^u\right) \mid_{u = \ln(x)}} \\ \pause
                          &= \dfrac{1}{e^{\ln(x)}} \pause = \dfrac{1}{x}
   \end{align} \pause

   \begin{alertblock}{Derivative of Natural Logarithm}\index{derivative!of natural logarithm}
     \begin{equation}
       \dfrac{d}{dx} \ln(x) = \dfrac{1}{x}
     \end{equation}
   \end{alertblock}
 \end{frame}
 
 \begin{frame}
   \frametitle{Inverse Differentiation IV}
   Ex. $\dfrac{d}{dx} \arcsin(x)$. \pause

   \begin{align}
     \dfrac{d}{dx} \arcsin(x) &= \dfrac{1}{\left(\dfrac{d}{du}\sin(u)\right)\mid_{u = \arcsin(x)}} \pause = \dfrac{1}{(\cos(u))\mid_{u = \arcsin(x)}} \\
                              &= \cos(\arcsin(x)) \pause = \dfrac{1}{\sqrt{1 - x^2}} \\ \pause
     \dfrac{d}{dx} \arccos(x) &= \dfrac{1}{-\sin(\arccos(x))} \pause = -\dfrac{1}{\sqrt{1 - x^2}} \\ \pause
     \dfrac{d}{dx} \arctan(x) &= \dfrac{1}{\sec^2(\arctan(x))} \pause = \dfrac{1}{1 + x^2}
   \end{align} 
 \end{frame}

 \begin{frame}
   \frametitle{Inverse Differentiation V}

   \begin{alertblock}{Derivatives of Inverse Trigonometric Functions}\index{derivative!of inverse trig functions}
     \begin{gather}
       \dfrac{d}{dx} \arcsin(x) = \dfrac{1}{\sqrt{1 - x^2}} \\
       \dfrac{d}{dx} \arccos(x) = -\dfrac{1}{\sqrt{1 - x^2}} \\
       \dfrac{d}{dx} \arctan(x) = \dfrac{1}{1 + x^2}
     \end{gather}
   \end{alertblock}
 \end{frame}

 \begin{frame}
   \frametitle{Implicit Differentiation I}
   What about for the other kinds of implicit relations? \pause

   Q: At a point $(x, y)$ on the circle of radius $1$, what is the slope of the line tangent to the circle at that point? \pause

   Circle of radius $1$ $\implies$ $x^2 + y^2 = 1$.

   \begin{gather}
      x^2 + y^2 = 1 \\
      \text{Differentiate both sides w.r.t. } x \\ \pause
      \dfrac{d}{dx}(x^2 + y^2) = \dfrac{d}{dx} 1 \\ \pause
      \dfrac{d}{dx}(x^2) + \dfrac{d}{dx}(y^2) = 0 \\ \pause
      2x + \dfrac{d}{dx} (y(x)^2) = 0 \\ \pause
      2x + 2y \cdot \dfrac{dy}{dx} = 0 \pause \quad 2y \dfrac{dy}{dx} = -2x \pause \quad \dfrac{dy}{dx} = -\dfrac{x}{y}
   \end{gather}
 \end{frame}

 \begin{frame}
   \frametitle{Implicit Differentiation II}

   The core idea is that in an implicit relation $F(x, y) = 0$,
   \pause
   we can \textit{always} treat one of the variables as a function of the other (fancy reason: \textit{implicit function theorem}) \pause
   and thus can differentiate implicit relations just by using the chain rule! \pause

    We'll see more application oriented examples when we do related rates soon :)
 \end{frame}

 \begin{frame}
   \frametitle{Worksheet Time}

   Let's get into groups and work on the worksheet!
 \end{frame}

 \begin{frame}
   \frametitle{End of Class}

   That's all folks!
 \end{frame}
\end{document}
